%% FedNPG-ADMM reproduction report
%%
%% SINGLE COLUMN, deliberately. Three reasons, all reversible:
%%   1. The paper we reproduce (NeurIPS 2023) is itself single-column, so a
%%      side-by-side reading of our tables against theirs stays honest.
%%   2. Our comparison tables are wide -- four agent counts x (ours, paper)
%%      -- and would need \resizebox or landscape in a two-column layout.
%%   3. Algorithm blocks and aligned equations survive line breaks better.
%% To try two columns: add [twocolumn] to \documentclass and expect to wrap
%% Tables 1-2 in table* and shrink the algorithm float.
\documentclass[11pt,a4paper]{article}

\usepackage[T1]{fontenc}
\usepackage[utf8]{inputenc}
\usepackage{lmodern}
\usepackage[margin=2.6cm]{geometry}

\usepackage{amsmath,amssymb,amsthm}
\usepackage{booktabs}
\usepackage{graphicx}
\usepackage{subcaption}
\usepackage[font=small,labelfont=bf]{caption}
\usepackage{algorithm}
\usepackage{algpseudocode}
\usepackage{siunitx}
\sisetup{detect-weight=true,detect-family=true}
% siunitx v2 (TeX Live <=2020) hides \byte behind binary-units; v3 provides
% it natively and removed the option. Apply it only on the old version.
\makeatletter
\@ifpackagelater{siunitx}{2021/01/01}{}{\sisetup{binary-units=true}}
\makeatother
\usepackage[dvipsnames]{xcolor}
\usepackage{pgfplots}
% TeX Live 2020 ships pgfplots with a 1.17 ceiling; 1.17 is also accepted
% by current pgfplots, so this compiles both here and on Overleaf.
\pgfplotsset{compat=1.17}
\usepgfplotslibrary{fillbetween}
% Figure palette. Colour-blind safe (Okabe-Ito), and distinguishable in
% greyscale by the line ordering rather than by hue alone.
\definecolor{colorA}{HTML}{0072B2}
\definecolor{colorB}{HTML}{D55E00}
\definecolor{colorC}{HTML}{009E73}
\definecolor{colorD}{HTML}{CC79A7}
\usepackage[numbers,sort&compress]{natbib}
\usepackage[hidelinks,breaklinks]{hyperref}
% The title contains \\ for visual line breaks, which hyperref cannot put into
% the PDF metadata string. Give it a flat version instead of letting it warn
% and guess.
\hypersetup{pdftitle={Communication-Efficient Federated Natural Policy
  Gradient: A Reproduction and Extension of FedNPG-ADMM},
  pdfauthor={Edoardo Scarpel, Daniele Zebele}}
\usepackage{microtype}
\usepackage{multirow}

% Notation shortcuts -- keep these consistent with Section 2.
\newcommand{\Rb}{\mathbb{R}}
\newcommand{\Exp}{\mathbb{E}}
\newcommand{\bg}{\mathbf{g}}
\newcommand{\by}{\mathbf{y}}
\newcommand{\bH}{\mathbf{H}}
\newcommand{\bF}{\mathbf{F}}
\newcommand{\bthe}{\boldsymbol{\theta}}
\newcommand{\blam}{\boldsymbol{\lambda}}
\newcommand{\Oh}{\mathcal{O}}
% Flag for anything not yet measured; delete the definition to find leftovers.
\newcommand{\TODO}[1]{\textcolor{BrickRed}{\textbf{[TODO: #1]}}}

\title{\textbf{Communication-Efficient Federated Natural \\
Policy Gradient: A Reproduction \\
and Extension of FedNPG-ADMM}}

\author{
  Edoardo Scarpel \\
  \texttt{edoardo.scarpel@phd.unipd.it}
  \and
  Daniele Zebele \\
  \texttt{daniele.zebele@phd.unipd.it}
}
\date{Department of Information Engineering, University of Padua \\[2pt] \today}

\begin{document}
\maketitle

\begin{abstract}
% !TEX root = ../main.tex
Federated Natural Policy Gradient (FedNPG) inherits the fast convergence of
Natural Policy Gradient (NPG) but
requires every agent to uplink its local Fisher information matrix $\bH_i \in
\Rb^{d\times d}$ each round. The resulting $\Oh(d^2)$ communication cost
becomes prohibitive for deep policies. \citet{lan2023fednpgadmm} reformulate
the global natural-gradient direction as a distributed quadratic program
solved by the Alternating Direction Method of Multipliers (ADMM), so that
each agent transmits only two $d$-dimensional vectors,
an $\Oh(d)$ cost, at the same convergence rate. We reproduce that method from
scratch in PyTorch and evaluate it on Swimmer-v4, Hopper-v4 and Humanoid-v4
with $N \in \{1,2,4,8\}$ agents over five random seeds. On Swimmer we match
the reported final rewards at every agent count, with tighter seed variance
than the original at each of them; on Humanoid the method scales from
$993$ to $6329$ average return as agents increase from one to eight. Hopper is
the one environment where we fail to reproduce the reported trend: performance
\emph{decreases} monotonically with the number of agents. We trace this to the
interaction between Hopper's small per-step survival bonus and the variance
reduction that additional agents provide, and support the explanation with
per-seed and per-step measurements. We further quantify the communication
saving directly, measuring reductions of $2.4\times 10^{3}$, $9.2\times
10^{3}$ and $3.6\times 10^{5}$ transmitted parameters per agent per round on
the three environments.
Finally, we quantize the agent-to-server transmissions, which the original
work does not attempt. With an unbiased stochastic quantizer at 8 bits the
uplink falls by a further $8\times$, compounding to $2.9\times 10^{6}$ fewer
transmitted bits per agent per round than standard FedNPG on Humanoid, while
final rewards stay within one standard deviation of the unquantized runs on all
three environments and agent counts. At 6 bits Swimmer and Hopper are still
unaffected, whereas Humanoid collapses.

\end{abstract}

% !TEX root = ../main.tex
\section{Introduction}
\label{sec:introduction}

Federated Reinforcement Learning (FedRL) lets several agents improve a shared
policy without pooling their raw trajectories. In the \emph{horizontal}
setting studied here, every agent faces the same environment, collects its own
experience, and periodically exchanges model information with a central server
\citep{qi2021fedrlsurvey}. This differs from Multi-Agent Reinforcement Learning (MARL), where agents
inhabit a single shared environment and interact with one another. No agent
here observes another's actions.

First-order federated policy gradient methods communicate a $d$-dimensional
gradient per round, which is cheap but inherits the slow convergence of plain
policy gradient. Natural Policy Gradient (NPG) converges considerably faster
by preconditioning the gradient with the inverse Fisher information matrix,
but a naive federation of NPG requires each agent to transmit its full local
Fisher $\bH_i \in \Rb^{d \times d}$. For the policies used in continuous
control this is fatal. Our Humanoid policy has $d = \num{727074}$ parameters,
so a single $\bH_i$ occupies \SI{2.1}{\peta\byte} at double precision, four
orders of magnitude beyond any machine we could deploy.

\citet{lan2023fednpgadmm} resolve this by observing that the global NPG direction
$\left(\sum_i \bH_i\right)^{-1}\sum_i \bg_i$ is the solution of a distributed
quadratic program, which consensus ADMM solves while exchanging only
$d$-dimensional vectors. The communication cost drops from $\Oh(d^2)$ to
$\Oh(d)$ at the same convergence rate and with no additional sample
complexity.

\paragraph{Contributions.} This report covers four deliverables.
\begin{enumerate}
  \item An exposition of the method and its communication argument
        (Sections~\ref{sec:background}--\ref{sec:method}).
  \item A from-scratch PyTorch implementation of Algorithm~\ref{alg:fednpgadmm}
        and a reproduction study across three MuJoCo environments, four agent
        counts and five seeds (Sections~\ref{sec:implementation}--%
        \ref{sec:reproduction}), including a negative result on Hopper that we
        diagnose rather than omit.
  \item A comparison against the two baselines the original work uses:
        standard-average FedNPG at $\Oh(d^2)$ and first-order FedPPO at
        $\Oh(d)$ (Section~\ref{sec:baselines}).
  \item Quantization of the agent-to-server transmissions, which is not part of
        the original paper and constitutes the novel contribution of this
        project (Section~\ref{sec:quantization}).
\end{enumerate}

Three of our results run against the original. The Hopper reproduction
contradicts the trend the authors report, our measured communication saving on
Swimmer is about half an order of magnitude smaller than the figure they quote,
and four of the hypotheses we formed about the Hopper discrepancy were refuted
by subsequent measurement. We treat each of these at length in
Section~\ref{sec:discussion}, since a reproduction is informative precisely
where it diverges, and since the refuted hypotheses constrain what the
remaining explanation can be.

% !TEX root = ../main.tex
\section{Background}
\label{sec:background}

\subsection{Setting}

$N$ agents share a policy $\pi_{\bthe}$ with parameters $\bthe \in \Rb^d$ and
each interacts with its own copy of the same Markov Decision Process (MDP) $(\mathcal{S}, \mathcal{A},
P, r, \gamma)$. The objective is the usual expected discounted return
\begin{equation}
  J(\bthe) \;=\; \Exp_{\tau \sim \pi_{\bthe}}
  \left[\textstyle\sum_{t=0}^{\infty} \gamma^{t} r(s_t, a_t)\right],
\end{equation}
and the federation seeks $\bthe^\star = \arg\max_{\bthe} J(\bthe)$ using every
agent's samples while transmitting as little as possible.

\subsection{Natural policy gradient}

Plain policy gradient ascends $\nabla_{\bthe} J$, which is sensitive to the
parameterization. An equal step in parameter space can mean a negligible or a
catastrophic change in the induced action distribution. NPG instead ascends
the gradient preconditioned by the inverse Fisher information matrix
\begin{equation}
  \bF(\bthe) \;=\;
  \Exp\!\left[\nabla_{\bthe} \log \pi_{\bthe}(a\mid s)\,
              \nabla_{\bthe} \log \pi_{\bthe}(a\mid s)^{\top}\right],
  \qquad
  \bthe \leftarrow \bthe + \alpha\, \bF(\bthe)^{-1} \nabla_{\bthe} J(\bthe),
\end{equation}
which makes the step invariant to reparameterization and yields a well-defined
trust region: constraining the average Kullback--Leibler (KL) divergence between the old and new
policies to $\delta$ gives the closed-form step size $\alpha = \sqrt{2\delta /
(\bg^{\top} \bF^{-1} \bg)}$ \citep{kakade2001npg, schulman2015trpo}.
Single-agent NPG, and Trust Region Policy Optimization (TRPO) with it, is
matrix-free, because $\bF$ is never needed
explicitly: only $\bF v$ for arbitrary $v$, obtainable by a double backward
pass through the KL divergence.

\subsection{The federated bottleneck}

Federating NPG naively means forming the global direction from every agent's
local gradient $\bg_i$ and local Fisher $\bH_i$:
\begin{equation}
  \bthe \leftarrow \bthe +
  \sqrt{\frac{2N\delta}
             {\big(\sum_{i} \bg_i\big)^{\!\top}
              \big(\sum_{i} \bH_i\big)^{-1} \sum_{i} \bg_i}}
  \left(\sum_{i=1}^{N} \bH_i\right)^{\!-1} \sum_{i=1}^{N} \bg_i .
  \label{eq:stdfednpg}
\end{equation}
This is what \citet{lan2023fednpgadmm} call \emph{standard-average FedNPG}.
The matrix-free trick no longer applies. The server needs each agent's $\bH_i$
to form the sum, and no agent can evaluate a Fisher-vector product on
another's data. Each agent must therefore uplink $\Oh(d^2)$ numbers per round.
Table~\ref{tab:commcost} shows what that means concretely for the three
environments studied here.

\begin{table}[t]
  \centering
  \caption{Per-agent uplink per iteration, at the paper's convention of 64 bits
  per transmitted parameter. $d$ is the policy parameter count for the
  architectures in Table~\ref{tab:table4}. The dense Fisher is reported for
  scale; it is never materialized in our implementation
  (Section~\ref{sec:baselines}).}
  \label{tab:commcost}
  \begin{tabular}{lrrrr}
    \toprule
    Environment & $d$ & Dense $\bH_i$ & $\Oh(d^2)$ / $\Oh(d)$ & Order \\
    \midrule
    Swimmer-v4  & \num{4868}   & \SI{0.095}{\giga\byte} & \num{2434}   & $10^{3.4}$ \\
    Hopper-v4   & \num{18438}  & \SI{1.36}{\giga\byte}  & \num{9219}   & $10^{4.0}$ \\
    Humanoid-v4 & \num{727074} & \SI{2114}{\giga\byte}  & \num{363537} & $10^{5.6}$ \\
    \bottomrule
  \end{tabular}
\end{table}

% !TEX root = ../main.tex
\section{FedNPG-ADMM}
\label{sec:method}

\subsection{The key reformulation}

The global direction in \eqref{eq:stdfednpg} is the stationary point of a
quadratic problem that decomposes across agents:
\begin{equation}
  \by^{\star}
  \;=\;\Big(\sum_i \bH_i\Big)^{\!-1}\sum_i \bg_i
  \;=\;\arg\max_{\by}\;
  \sum_{i=1}^{N}\Big(\bg_i^{\top}\by - \tfrac{1}{2}\by^{\top}\bH_i\by\Big).
  \label{eq:qp}
\end{equation}
Each summand depends only on agent $i$'s own data. Introducing per-agent
copies $\by_i$ with the consensus constraint $\by_i = \by$ for all $i$ turns
\eqref{eq:qp} into a form consensus ADMM solves by alternating a local solve, a
dual update and an average. The only quantities crossing the network are
$d$-dimensional.

\subsection{Algorithm}

\begin{algorithm}[t]
\caption{FedNPG-ADMM (Algorithm 1 of \citealp{lan2023fednpgadmm})}
\label{alg:fednpgadmm}
\begin{algorithmic}[1]
  \State \textbf{Initialize} $\bthe^0$, $\by^0 = \mathbf{0}$,
         $\by_i^0 = \by^0$, $\blam_i = \mathbf{0}$ for all $i$
  \For{$k = 0, 1, \dots, K-1$}
    \State Server broadcasts $\bthe^k$ and $\by^{k-1}$ to all agents
    \For{each agent $i$ \textbf{in parallel}}
      \State Collect a rollout of $T$ timesteps under $\pi_{\bthe^k}$
      \State $\blam_i \leftarrow \blam_i + \rho\,(\by_i - \by^{k-1})$
             \Comment{dual update}
      \State $\bg_i \leftarrow \frac{1}{|\mathcal{D}_i|}\sum_{t}
             \nabla_{\bthe}\log\pi_{\bthe}(a_t \mid s_t)\,\hat{A}_t$
             \Comment{score-function gradient}
      \State $\by_i \leftarrow (\bH_i + \rho I)^{-1}
             (\bg_i - \blam_i + \rho\,\by^{k-1})$
             \Comment{local solve, by conjugate gradient}
      \State Send $\by_i, \bg_i$ to the server
             \Comment{$\Oh(d)$}
    \EndFor
    \State $\by^{k} \leftarrow \frac{1}{N}\sum_{i} \by_i$
    \State $\bthe^{k+1} \leftarrow \bthe^{k} +
           \eta\sqrt{\dfrac{2N\delta}{(\sum_i \bg_i)^{\top}\by^{k}}}\;\by^{k}$
  \EndFor
\end{algorithmic}
\end{algorithm}

Line~8 is where the $\Oh(d)$ claim is earned. The local system $(\bH_i + \rho
I)\by_i = \bg_i - \blam_i + \rho\by^{k-1}$ is solved by Conjugate Gradient (CG),
which needs only Fisher-vector products on agent $i$'s own data, so $\bH_i$ is
never formed. The $\rho I$ term is simultaneously the ADMM penalty and a
Levenberg--Marquardt damping that keeps the system well-conditioned, so $\rho$
plays two roles at once.

\subsection{Why the step size carries \texorpdfstring{$N$}{N}}

At exact consensus, $\by^k$ solves \eqref{eq:qp}, hence $(\sum_i
\bg_i)^{\top}\by = \by^{\top}(\sum_i \bH_i)\by$. Constraining the
\emph{average} per-agent KL to $\delta$ means
$\tfrac12 \alpha^2 \by^{\top}\big(\tfrac{1}{N}\sum_i\bH_i\big)\by \le \delta$,
which rearranges exactly to the $\sqrt{2N\delta/(\sum_i\bg_i)^\top\by}$ of
line~12. The factor $N$ is therefore not a free scaling: it is what makes the
trust region a statement about the average agent rather than the sum.

\subsection{A property of the single-agent case}

One consequence of Algorithm~\ref{alg:fednpgadmm} matters for interpreting
$N=1$ results and is not remarked on in the original. With one agent, the
broadcast consensus variable is that agent's own previous solve, $\by^{k-1} =
\by_1^{k-1}$, so line~6 reads $\blam_1 \leftarrow \blam_1 + \rho(\by_1 -
\by_1) = \blam_1$ and the dual variable never leaves its zero initialization.
Algorithm~\ref{alg:fednpgadmm} at $N = 1$ therefore reduces exactly to
\begin{equation}
  \by \;=\; (\bH + \rho I)^{-1}\big(\bg + \rho\,\by^{k-1}\big),
\end{equation}
damped NPG with a momentum term, exercising none of the consensus machinery.
We verified this numerically: $\|\blam_1\| = 0$ to machine precision at every
iteration for $N=1$, and non-zero from the second iteration onwards for $N >
1$. Any $N = 1$ number in this report, and in the original, should be read
with that in mind.

% !TEX root = ../main.tex
\section{Implementation}
\label{sec:implementation}

The implementation is written from scratch in PyTorch; no reference code for
FedNPG-ADMM is publicly available. The algorithm proper (\texttt{agent.py},
\texttt{server.py}) stays free of orchestration concerns, and every
agent-to-server transmission is routed through a single \texttt{Channel}
abstraction that meters bytes. That chokepoint is what makes the communication
comparison in Section~\ref{sec:baselines} meaningful: all three algorithms
count bytes through the same code, and the quantization study of
Section~\ref{sec:quantization} attaches there without any algorithm importing
it.

\subsection{What comes from the paper and what does not}

Table~\ref{tab:table4} reproduces the paper's hyperparameters. Two values in
our pipeline are \emph{not} transcriptions. First, $\eta$ in line~12 of
Algorithm~\ref{alg:fednpgadmm} appears nowhere in Table~4 of the original. The
paper derives it analytically from smoothness constants it never reports
numerically, so we set $\eta = 1$, taking the ADMM step unscaled. Second,
Table~4's ``Learning rate'' column ($3\times10^{-4}$, $3\times10^{-4}$,
$10^{-5}$) is used here as the \emph{critic's} Adam learning rate, not an
actor step size. Those are characteristic Adam values, and the paper defines
the actor's step analytically, which we read as evidence that the column
refers to the critic. This is an interpretation of an ambiguous table.

\subsection{Practical settings}

The paper's own setup section states that it follows ``the practical settings
in TRPO with line search and in stable-baselines3 with generalized advantage
estimation (0.95) and the Adam optimizer''. Algorithm~\ref{alg:fednpgadmm}'s
pseudocode contains none of these, so we adopt the standard versions and list
them in Appendix~\ref{app:hyperparameters}: running observation and reward
normalization, advantage normalization, orthogonal initialization, Generalized Advantage
Estimation (GAE) with $\lambda = 0.95$, minibatch critic fitting, and a backtracking line search on
the true KL and the importance-sampled surrogate. The line search matters. The
closed-form step of line~12 controls only the second-order approximation of
the KL, so without a safeguard it can overshoot the intended trust region.

Two further design choices are ours because the paper is silent. The critic is
fit on data \emph{pooled} across agents; the paper describes only the actor's
update and never says how a federated critic is trained. And when $(\sum_i
\bg_i)^{\top}\by \le 0$ the step size in line~12 is undefined, in which case
we skip the update for that iteration and log how often it happens.

\subsection{Reporting conventions}

We report the paper's own metric: the average trajectory return across the
batch collected at each iteration, plotted against iterations. Two measurement
issues are worth stating because they materially changed our numbers.

First, this metric averages only the episodes that \emph{complete} within one
iteration's timestep budget. On Hopper at $N=1$ that is roughly $1.2$ episodes
per point, so each plotted value is close to a single episode, which makes the
curve appear unstable even when the underlying policy is not. We therefore
also record a trailing 100-episode mean and, separately, a deterministic
evaluation (mean action, no exploration noise, on a held-out environment), the
quantity published benchmarks such as \texttt{rl-baselines3-zoo} report. The
deterministic number runs roughly $1.2\times$ higher than the stochastic one,
so the two are not interchangeable.

Second, no published MuJoCo learning curve plots this metric raw.
\citet{spinningup} smooth the identical quantity over an 11-epoch window and
average 10 seeds, and \texttt{rl-baselines3-zoo} uses a 100-episode rolling
window. Where we show curves we state the smoothing in the caption, following
the recommendation of \citet{huang2024openrlbenchmark}.

% !TEX root = ../main.tex
\section{Reproduction results}
\label{sec:reproduction}

We run $N \in \{1,2,4,8\}$ on Swimmer-v4, Hopper-v4 and Humanoid-v4 with five
random seeds per configuration (60 runs). Table~\ref{tab:reproduction} reports
the mean of each seed's final 100 iterations, then the mean and standard
deviation across seeds, the same quantity Table~2 of the original reports.

\begin{table}[t]
  \centering
  \caption{Final rewards, mean $\pm$ standard deviation across five seeds.
  Original values are from Table~2 of \citet{lan2023fednpgadmm}, computed over
  ten seeds. Humanoid does not appear in their Table~2.}
  \label{tab:reproduction}
  \begin{tabular}{llrrrr}
    \toprule
    Environment & Source & $N=1$ & $N=2$ & $N=4$ & $N=8$ \\
    \midrule
    \multirow{2}{*}{Swimmer-v4}
      & Ours  & $\mathbf{123.9 \pm 2.3}$ & $126.0 \pm 1.3$ & $127.0 \pm 0.8$ & $\mathbf{128.4 \pm 0.8}$ \\
      & Paper & $109.4 \pm 5.9$  & $123.6 \pm 10.3$ & $127.2 \pm 2.2$ & $128.5 \pm 1.9$ \\
    \midrule
    \multirow{2}{*}{Hopper-v4}
      & Ours  & $\mathbf{2497 \pm 635}$ & $2140 \pm 386$ & $1658 \pm 273$ & $1605 \pm 159$ \\
      & Paper & $1473 \pm 383$ & $2384 \pm 371$ & $2507 \pm 230$ & $2719 \pm 173$ \\
    \midrule
    Humanoid-v4
      & Ours  & $993 \pm 187$ & $3391 \pm 796$ & $5343 \pm 121$ & $\mathbf{6329 \pm 334}$ \\
    \bottomrule
  \end{tabular}
\end{table}

\begin{figure}[t]
  \centering
  \begin{subfigure}[b]{0.49\textwidth}
    %% Convergence, Swimmer-v4. Mean over five seeds; band spans min to max.
%% Data: figures/convergence_swimmer.dat (scripts/export_figures.py).
\begin{tikzpicture}
\begin{axis}[
    width=\linewidth,
    height=0.78\linewidth,
    xlabel={Iteration},
    ylabel={Average batch return},
    xmin=1, xmax=1000,
    ymin=0, ymax=140,
    xtick distance=250,
    ytick distance=35,
    ymajorgrids=true, xmajorgrids=true, grid style=dashed,
    tick label style={font=\footnotesize},
    label style={font=\footnotesize},
    legend style={legend columns=2, font=\footnotesize, cells={anchor=west},
                  anchor=south east, at={(0.99,0.02)}},
]

\addplot[name path=N1lo, draw=none, forget plot]
    table[x=iteration, y=N1_lo] {figures/convergence_swimmer.dat};
\addplot[name path=N1hi, draw=none, forget plot]
    table[x=iteration, y=N1_hi] {figures/convergence_swimmer.dat};
\addplot[colorA, opacity=0.15, forget plot] fill between[of=N1lo and N1hi];
\addplot[colorA, line width=0.9pt]
    table[x=iteration, y=N1_mean] {figures/convergence_swimmer.dat};
\addlegendentry{$N=1$};

\addplot[name path=N2lo, draw=none, forget plot]
    table[x=iteration, y=N2_lo] {figures/convergence_swimmer.dat};
\addplot[name path=N2hi, draw=none, forget plot]
    table[x=iteration, y=N2_hi] {figures/convergence_swimmer.dat};
\addplot[colorB, opacity=0.15, forget plot] fill between[of=N2lo and N2hi];
\addplot[colorB, line width=0.9pt]
    table[x=iteration, y=N2_mean] {figures/convergence_swimmer.dat};
\addlegendentry{$N=2$};

\addplot[name path=N4lo, draw=none, forget plot]
    table[x=iteration, y=N4_lo] {figures/convergence_swimmer.dat};
\addplot[name path=N4hi, draw=none, forget plot]
    table[x=iteration, y=N4_hi] {figures/convergence_swimmer.dat};
\addplot[colorC, opacity=0.15, forget plot] fill between[of=N4lo and N4hi];
\addplot[colorC, line width=0.9pt]
    table[x=iteration, y=N4_mean] {figures/convergence_swimmer.dat};
\addlegendentry{$N=4$};

\addplot[name path=N8lo, draw=none, forget plot]
    table[x=iteration, y=N8_lo] {figures/convergence_swimmer.dat};
\addplot[name path=N8hi, draw=none, forget plot]
    table[x=iteration, y=N8_hi] {figures/convergence_swimmer.dat};
\addplot[colorD, opacity=0.15, forget plot] fill between[of=N8lo and N8hi];
\addplot[colorD, line width=0.9pt]
    table[x=iteration, y=N8_mean] {figures/convergence_swimmer.dat};
\addlegendentry{$N=8$};
\end{axis}
\end{tikzpicture}

    \caption{Swimmer-v4}
  \end{subfigure}
  \hfill
  \begin{subfigure}[b]{0.49\textwidth}
    %% Convergence, Hopper-v4. Mean over five seeds; band spans min to max.
%% Data: figures/convergence_hopper.dat (scripts/export_figures.py).
\begin{tikzpicture}
\begin{axis}[
    width=\linewidth,
    height=0.78\linewidth,
    xlabel={Iteration},
    ylabel={Average batch return},
    xmin=1, xmax=2000,
    ymin=0, ymax=3600,
    xtick distance=500,
    ytick distance=900,
    scaled y ticks=base 10:-3,
    ymajorgrids=true, xmajorgrids=true, grid style=dashed,
    tick label style={font=\footnotesize},
    label style={font=\footnotesize},
    legend style={legend columns=2, font=\footnotesize, cells={anchor=west},
                  anchor=south east, at={(0.99,0.02)}},
]

\addplot[name path=N1lo, draw=none, forget plot]
    table[x=iteration, y=N1_lo] {figures/convergence_hopper.dat};
\addplot[name path=N1hi, draw=none, forget plot]
    table[x=iteration, y=N1_hi] {figures/convergence_hopper.dat};
\addplot[colorA, opacity=0.15, forget plot] fill between[of=N1lo and N1hi];
\addplot[colorA, line width=0.9pt]
    table[x=iteration, y=N1_mean] {figures/convergence_hopper.dat};
\addlegendentry{$N=1$};

\addplot[name path=N2lo, draw=none, forget plot]
    table[x=iteration, y=N2_lo] {figures/convergence_hopper.dat};
\addplot[name path=N2hi, draw=none, forget plot]
    table[x=iteration, y=N2_hi] {figures/convergence_hopper.dat};
\addplot[colorB, opacity=0.15, forget plot] fill between[of=N2lo and N2hi];
\addplot[colorB, line width=0.9pt]
    table[x=iteration, y=N2_mean] {figures/convergence_hopper.dat};
\addlegendentry{$N=2$};

\addplot[name path=N4lo, draw=none, forget plot]
    table[x=iteration, y=N4_lo] {figures/convergence_hopper.dat};
\addplot[name path=N4hi, draw=none, forget plot]
    table[x=iteration, y=N4_hi] {figures/convergence_hopper.dat};
\addplot[colorC, opacity=0.15, forget plot] fill between[of=N4lo and N4hi];
\addplot[colorC, line width=0.9pt]
    table[x=iteration, y=N4_mean] {figures/convergence_hopper.dat};
\addlegendentry{$N=4$};

\addplot[name path=N8lo, draw=none, forget plot]
    table[x=iteration, y=N8_lo] {figures/convergence_hopper.dat};
\addplot[name path=N8hi, draw=none, forget plot]
    table[x=iteration, y=N8_hi] {figures/convergence_hopper.dat};
\addplot[colorD, opacity=0.15, forget plot] fill between[of=N8lo and N8hi];
\addplot[colorD, line width=0.9pt]
    table[x=iteration, y=N8_mean] {figures/convergence_hopper.dat};
\addlegendentry{$N=8$};
\end{axis}
\end{tikzpicture}

    \caption{Hopper-v4}
  \end{subfigure}

  \vspace{0.6em}
  \begin{subfigure}[b]{0.49\textwidth}
    %% Convergence, Humanoid-v4. Mean over five seeds; band spans min to max.
%% Data: figures/convergence_humanoid.dat (scripts/export_figures.py).
\begin{tikzpicture}
\begin{axis}[
    width=\linewidth,
    height=0.78\linewidth,
    xlabel={Iteration},
    ylabel={Average batch return},
    xmin=1, xmax=3000,
    ymin=0, ymax=7000,
    xtick distance=1000,
    ytick distance=1750,
    scaled y ticks=base 10:-3,
    ymajorgrids=true, xmajorgrids=true, grid style=dashed,
    tick label style={font=\footnotesize},
    label style={font=\footnotesize},
    legend style={legend columns=2, font=\footnotesize, cells={anchor=west},
                  anchor=north west, at={(0.02,0.98)}},
]

\addplot[name path=N1lo, draw=none, forget plot]
    table[x=iteration, y=N1_lo] {figures/convergence_humanoid.dat};
\addplot[name path=N1hi, draw=none, forget plot]
    table[x=iteration, y=N1_hi] {figures/convergence_humanoid.dat};
\addplot[colorA, opacity=0.15, forget plot] fill between[of=N1lo and N1hi];
\addplot[colorA, line width=0.9pt]
    table[x=iteration, y=N1_mean] {figures/convergence_humanoid.dat};
\addlegendentry{$N=1$};

\addplot[name path=N2lo, draw=none, forget plot]
    table[x=iteration, y=N2_lo] {figures/convergence_humanoid.dat};
\addplot[name path=N2hi, draw=none, forget plot]
    table[x=iteration, y=N2_hi] {figures/convergence_humanoid.dat};
\addplot[colorB, opacity=0.15, forget plot] fill between[of=N2lo and N2hi];
\addplot[colorB, line width=0.9pt]
    table[x=iteration, y=N2_mean] {figures/convergence_humanoid.dat};
\addlegendentry{$N=2$};

\addplot[name path=N4lo, draw=none, forget plot]
    table[x=iteration, y=N4_lo] {figures/convergence_humanoid.dat};
\addplot[name path=N4hi, draw=none, forget plot]
    table[x=iteration, y=N4_hi] {figures/convergence_humanoid.dat};
\addplot[colorC, opacity=0.15, forget plot] fill between[of=N4lo and N4hi];
\addplot[colorC, line width=0.9pt]
    table[x=iteration, y=N4_mean] {figures/convergence_humanoid.dat};
\addlegendentry{$N=4$};

\addplot[name path=N8lo, draw=none, forget plot]
    table[x=iteration, y=N8_lo] {figures/convergence_humanoid.dat};
\addplot[name path=N8hi, draw=none, forget plot]
    table[x=iteration, y=N8_hi] {figures/convergence_humanoid.dat};
\addplot[colorD, opacity=0.15, forget plot] fill between[of=N8lo and N8hi];
\addplot[colorD, line width=0.9pt]
    table[x=iteration, y=N8_mean] {figures/convergence_humanoid.dat};
\addlegendentry{$N=8$};
\end{axis}
\end{tikzpicture}

    \caption{Humanoid-v4}
  \end{subfigure}
  \caption{Convergence for each agent count. The solid line is the mean over five
  seeds and the shaded band spans the minimum and maximum seed, so the band
  shows the full observed range and not a modelled interval. Curves are
  smoothed with a centred 21-iteration rolling mean before aggregation,
  following \citet{spinningup}, who smooth the same quantity over eleven
  epochs. Swimmer and Humanoid separate cleanly by agent count. On Hopper the
  bands overlap throughout and the ordering reverses, which is the visual form
  of the result discussed in Section~\ref{sec:hopper}.}
  \label{fig:convergence}
\end{figure}

\begin{figure}[t]
  \centering
  \begin{subfigure}[b]{0.49\textwidth}
    %% Return against cumulative uplink traffic, Swimmer-v4.
%% x is the federation total: the per-agent cost is 2d*8 bytes per iteration
%% for every N, and that invariance is the O(d) claim itself.
%% Data: figures/communication_swimmer.dat (scripts/export_figures.py).
\begin{tikzpicture}
\begin{axis}[
    width=\linewidth,
    height=0.78\linewidth,
    xlabel={Cumulative uplink (bytes)},
    ylabel={Average batch return},
    xmode=log,
    ymin=0, ymax=140,
    ytick distance=35,
    ymajorgrids=true, xmajorgrids=true, grid style=dashed,
    tick label style={font=\footnotesize},
    label style={font=\footnotesize},
    legend style={legend columns=2, font=\footnotesize, cells={anchor=west},
                  anchor=north west, at={(0.02,0.98)}},
]

\addplot[name path=N1clo, draw=none, forget plot]
    table[x=N1_bytes, y=N1_lo] {figures/communication_swimmer.dat};
\addplot[name path=N1chi, draw=none, forget plot]
    table[x=N1_bytes, y=N1_hi] {figures/communication_swimmer.dat};
\addplot[colorA, opacity=0.15, forget plot] fill between[of=N1clo and N1chi];
\addplot[colorA, line width=0.9pt]
    table[x=N1_bytes, y=N1_mean] {figures/communication_swimmer.dat};
\addlegendentry{$N=1$};

\addplot[name path=N2clo, draw=none, forget plot]
    table[x=N2_bytes, y=N2_lo] {figures/communication_swimmer.dat};
\addplot[name path=N2chi, draw=none, forget plot]
    table[x=N2_bytes, y=N2_hi] {figures/communication_swimmer.dat};
\addplot[colorB, opacity=0.15, forget plot] fill between[of=N2clo and N2chi];
\addplot[colorB, line width=0.9pt]
    table[x=N2_bytes, y=N2_mean] {figures/communication_swimmer.dat};
\addlegendentry{$N=2$};

\addplot[name path=N4clo, draw=none, forget plot]
    table[x=N4_bytes, y=N4_lo] {figures/communication_swimmer.dat};
\addplot[name path=N4chi, draw=none, forget plot]
    table[x=N4_bytes, y=N4_hi] {figures/communication_swimmer.dat};
\addplot[colorC, opacity=0.15, forget plot] fill between[of=N4clo and N4chi];
\addplot[colorC, line width=0.9pt]
    table[x=N4_bytes, y=N4_mean] {figures/communication_swimmer.dat};
\addlegendentry{$N=4$};

\addplot[name path=N8clo, draw=none, forget plot]
    table[x=N8_bytes, y=N8_lo] {figures/communication_swimmer.dat};
\addplot[name path=N8chi, draw=none, forget plot]
    table[x=N8_bytes, y=N8_hi] {figures/communication_swimmer.dat};
\addplot[colorD, opacity=0.15, forget plot] fill between[of=N8clo and N8chi];
\addplot[colorD, line width=0.9pt]
    table[x=N8_bytes, y=N8_mean] {figures/communication_swimmer.dat};
\addlegendentry{$N=8$};
\end{axis}
\end{tikzpicture}

    \caption{Swimmer-v4}
  \end{subfigure}
  \hfill
  \begin{subfigure}[b]{0.49\textwidth}
    %% Return against cumulative uplink traffic, Hopper-v4.
%% x is the federation total: the per-agent cost is 2d*8 bytes per iteration
%% for every N, and that invariance is the O(d) claim itself.
%% Data: figures/communication_hopper.dat (scripts/export_figures.py).
\begin{tikzpicture}
\begin{axis}[
    width=\linewidth,
    height=0.78\linewidth,
    xlabel={Cumulative uplink (bytes)},
    ylabel={Average batch return},
    xmode=log,
    ymin=0, ymax=3600,
    ytick distance=900,
    scaled y ticks=base 10:-3,
    ymajorgrids=true, xmajorgrids=true, grid style=dashed,
    tick label style={font=\footnotesize},
    label style={font=\footnotesize},
    legend style={legend columns=2, font=\footnotesize, cells={anchor=west},
                  anchor=north west, at={(0.02,0.98)}},
]

\addplot[name path=N1clo, draw=none, forget plot]
    table[x=N1_bytes, y=N1_lo] {figures/communication_hopper.dat};
\addplot[name path=N1chi, draw=none, forget plot]
    table[x=N1_bytes, y=N1_hi] {figures/communication_hopper.dat};
\addplot[colorA, opacity=0.15, forget plot] fill between[of=N1clo and N1chi];
\addplot[colorA, line width=0.9pt]
    table[x=N1_bytes, y=N1_mean] {figures/communication_hopper.dat};
\addlegendentry{$N=1$};

\addplot[name path=N2clo, draw=none, forget plot]
    table[x=N2_bytes, y=N2_lo] {figures/communication_hopper.dat};
\addplot[name path=N2chi, draw=none, forget plot]
    table[x=N2_bytes, y=N2_hi] {figures/communication_hopper.dat};
\addplot[colorB, opacity=0.15, forget plot] fill between[of=N2clo and N2chi];
\addplot[colorB, line width=0.9pt]
    table[x=N2_bytes, y=N2_mean] {figures/communication_hopper.dat};
\addlegendentry{$N=2$};

\addplot[name path=N4clo, draw=none, forget plot]
    table[x=N4_bytes, y=N4_lo] {figures/communication_hopper.dat};
\addplot[name path=N4chi, draw=none, forget plot]
    table[x=N4_bytes, y=N4_hi] {figures/communication_hopper.dat};
\addplot[colorC, opacity=0.15, forget plot] fill between[of=N4clo and N4chi];
\addplot[colorC, line width=0.9pt]
    table[x=N4_bytes, y=N4_mean] {figures/communication_hopper.dat};
\addlegendentry{$N=4$};

\addplot[name path=N8clo, draw=none, forget plot]
    table[x=N8_bytes, y=N8_lo] {figures/communication_hopper.dat};
\addplot[name path=N8chi, draw=none, forget plot]
    table[x=N8_bytes, y=N8_hi] {figures/communication_hopper.dat};
\addplot[colorD, opacity=0.15, forget plot] fill between[of=N8clo and N8chi];
\addplot[colorD, line width=0.9pt]
    table[x=N8_bytes, y=N8_mean] {figures/communication_hopper.dat};
\addlegendentry{$N=8$};
\end{axis}
\end{tikzpicture}

    \caption{Hopper-v4}
  \end{subfigure}

  \vspace{0.6em}
  \begin{subfigure}[b]{0.49\textwidth}
    %% Return against cumulative uplink traffic, Humanoid-v4.
%% x is the federation total: the per-agent cost is 2d*8 bytes per iteration
%% for every N, and that invariance is the O(d) claim itself.
%% Data: figures/communication_humanoid.dat (scripts/export_figures.py).
\begin{tikzpicture}
\begin{axis}[
    width=\linewidth,
    height=0.78\linewidth,
    xlabel={Cumulative uplink (bytes)},
    ylabel={Average batch return},
    xmode=log,
    ymin=0, ymax=7000,
    ytick distance=1750,
    scaled y ticks=base 10:-3,
    ymajorgrids=true, xmajorgrids=true, grid style=dashed,
    tick label style={font=\footnotesize},
    label style={font=\footnotesize},
    legend style={legend columns=2, font=\footnotesize, cells={anchor=west},
                  anchor=north west, at={(0.02,0.98)}},
]

\addplot[name path=N1clo, draw=none, forget plot]
    table[x=N1_bytes, y=N1_lo] {figures/communication_humanoid.dat};
\addplot[name path=N1chi, draw=none, forget plot]
    table[x=N1_bytes, y=N1_hi] {figures/communication_humanoid.dat};
\addplot[colorA, opacity=0.15, forget plot] fill between[of=N1clo and N1chi];
\addplot[colorA, line width=0.9pt]
    table[x=N1_bytes, y=N1_mean] {figures/communication_humanoid.dat};
\addlegendentry{$N=1$};

\addplot[name path=N2clo, draw=none, forget plot]
    table[x=N2_bytes, y=N2_lo] {figures/communication_humanoid.dat};
\addplot[name path=N2chi, draw=none, forget plot]
    table[x=N2_bytes, y=N2_hi] {figures/communication_humanoid.dat};
\addplot[colorB, opacity=0.15, forget plot] fill between[of=N2clo and N2chi];
\addplot[colorB, line width=0.9pt]
    table[x=N2_bytes, y=N2_mean] {figures/communication_humanoid.dat};
\addlegendentry{$N=2$};

\addplot[name path=N4clo, draw=none, forget plot]
    table[x=N4_bytes, y=N4_lo] {figures/communication_humanoid.dat};
\addplot[name path=N4chi, draw=none, forget plot]
    table[x=N4_bytes, y=N4_hi] {figures/communication_humanoid.dat};
\addplot[colorC, opacity=0.15, forget plot] fill between[of=N4clo and N4chi];
\addplot[colorC, line width=0.9pt]
    table[x=N4_bytes, y=N4_mean] {figures/communication_humanoid.dat};
\addlegendentry{$N=4$};

\addplot[name path=N8clo, draw=none, forget plot]
    table[x=N8_bytes, y=N8_lo] {figures/communication_humanoid.dat};
\addplot[name path=N8chi, draw=none, forget plot]
    table[x=N8_bytes, y=N8_hi] {figures/communication_humanoid.dat};
\addplot[colorD, opacity=0.15, forget plot] fill between[of=N8clo and N8chi];
\addplot[colorD, line width=0.9pt]
    table[x=N8_bytes, y=N8_mean] {figures/communication_humanoid.dat};
\addlegendentry{$N=8$};
\end{axis}
\end{tikzpicture}

    \caption{Humanoid-v4}
  \end{subfigure}
  \caption{The same runs plotted against cumulative uplink traffic instead of
  iteration count, on a logarithmic axis. The horizontal axis is the total
  across the federation: each agent transmits $2d$ values per iteration
  whatever $N$ is, so a per-agent axis would place all four series on top of
  one another, and that invariance is the $\Oh(d)$ property itself. Read this
  way the figure prices the extra agents, and the curves shift right as $N$
  grows.}
  \label{fig:communication}
\end{figure}

\subsection{Swimmer: reproduced}

Figure~\ref{fig:convergence}(a) shows the curves behind the Swimmer column, and
Figure~\ref{fig:communication}(a) the same runs against uplink traffic.
Our numbers fall within the reported confidence intervals at every agent
count, they are monotone in $N$ as the paper describes, and their seed variance
is in fact \emph{tighter} than the original's at every count, by a factor of
roughly two at both ends ($\pm 2.3$ against $\pm 5.9$ at $N=1$, $\pm 0.8$
against $\pm 1.9$ at $N=8$). At $N=1$ we also exceed the reported mean, $123.9$
against $109.4$. On this environment we therefore regard the result as
reproduced without qualification.

The variance does still fall with $N$, from $\pm 2.3$ at one agent to $\pm 0.8$
at eight, which is the direction the paper reports. The effect is modest here
because every seed converges: none of the five single-agent runs falls into the
local optimum Swimmer is known for, so there is no failure mode for the extra
agents to suppress. Humanoid in Section~\ref{sec:humanoid} shows the same claim
with far more room to act.

\subsection{Humanoid: strong scaling}
\label{sec:humanoid}

Humanoid is absent from the original Table~2, so we have no direct comparison.
The qualitative claim nevertheless holds, since return grows $6.4\times$ from $N=1$ to $N=8$,
monotonically, with the $N=4$ and $N=8$ variance small relative to the mean.

\subsection{Hopper: not reproduced}
\label{sec:hopper}

Hopper is the one environment where our result contradicts the original. Our
$N=1$ \emph{exceeds} the reported value ($2497$ against $1473$), but
performance then \emph{decreases} monotonically with more agents, ending at
$1605$ where the paper reports $2719$. The trend is inverted.

Section~\ref{sec:discussion} lists the aggregation hypotheses we tested and
refuted. The explanation consistent with the evidence is a
speed-versus-survival trade-off specific to Hopper's reward structure.
Table~\ref{tab:hopper} shows that the runs with more agents earn \emph{more}
reward per timestep while surviving \emph{fewer} timesteps: they converge to a
faster, less stable gait.

\begin{table}[t]
  \centering
  \caption{Hopper-v4: reward per step against episode length, final 100
  iterations. More agents yield a faster gait that falls sooner.}
  \label{tab:hopper}
  \begin{tabular}{lrrr}
    \toprule
    Configuration & Reward/step & Episode length & Return \\
    \midrule
    $N=1$ (best seed) & $3.76$ & $779$ & $2928$ \\
    $N=2$ & $4.20$ & $372$ & $1562$ \\
    $N=4$ & $4.27$ & $384$ & $1638$ \\
    $N=8$ & $\mathbf{4.49}$ & $\mathbf{368}$ & $1652$ \\
    \bottomrule
  \end{tabular}
\end{table}

Hopper's \texttt{healthy\_reward} is $1.0$ per surviving step, small relative
to its forward-velocity reward, so a gait that runs fast and falls at step
$\sim$400 is a local optimum. Humanoid's \texttt{healthy\_reward} is $5.0$,
which aligns survival with scoring, and there reward per step and episode
length rise
\emph{together} with $N$ ($5.70 \to 7.09$ and $206 \to 874$). The mechanism we propose is that additional agents reduce gradient noise, and
with it the implicit exploration that noise supplies. Whether the loss of that
exploration helps or hurts then depends on which basin the now more decisive
optimizer settles into. On Swimmer and Humanoid it steers away from the poor
optimum, and the seed variance falls accordingly; on Hopper it steers into one.
The same property of the method therefore produces the paper's result on two
environments and the opposite on the third.

\subsection{On the variance of the reported metric}
\label{sec:variance}

Single-seed Hopper and Humanoid curves swing between roughly $1000$ and $3700$
from one iteration to the next, while Swimmer's are flat to within $2\%$. This
is not instability of the algorithm. Swimmer-v4 has no termination condition
at all, and \citet{gymnasium} state that ``the Swimmer never terminates'', so
every episode runs exactly $1000$ steps. Hopper and Humanoid terminate when
the robot falls, which makes return approximately affine in survival time and
the per-batch mean a draw from a two-component mixture: roughly $58\%$ of
episodes reach the time limit near $3700$, while $9\%$ fall early below
$2000$. With one to two episodes per plotted point, the swing is the expected
behaviour of the estimator. The same policy evaluated deterministically scores
$3749 \pm 4$ and almost never falls.

We verified that the swing cannot be tuned away. Removing the
exploration-noise floor, removing the ceiling, and annealing the noise to a
level at which a fixed policy provably never falls all left the fall rate at
$9$--$12\%$. A policy trained at lower noise re-optimizes to a correspondingly
more marginal gait, so the failure rate is an equilibrium of the objective
rather than a function of the noise level. A never-falling policy would score
lower in expectation.

% !TEX root = ../main.tex
\section{Baseline comparison}
\label{sec:baselines}

The original work compares three methods at $N = 8$: standard-average FedNPG
\eqref{eq:stdfednpg}, FedNPG-ADMM, and first-order federated Proximal Policy Optimization (FedPPO) with clipping
parameter $0.2$. We implement all three in the same codebase, holding
everything except the server-side aggregation identical: same environments,
normalization, GAE, pooled critic and logged metrics. The comparison therefore
measures the algorithms, not the surrounding machinery.

\subsection{Standard-average FedNPG}

This arm solves \eqref{eq:stdfednpg} exactly by conjugate gradient on the
pooled Fisher, exploiting $\sum_i \bH_i v = N \cdot \bF_{\text{pooled}} v$.
The solve is matrix-free, but the \emph{metering} is not: the channel is
charged the full $d^2$ elements per agent per round.

The obvious objection is that a baseline meant to pay an $\Oh(d^2)$ cost ought
to pay it. Materializing $\sum_i \bH_i$ is impossible at Humanoid scale
(Table~\ref{tab:commcost}: \SI{2.1}{\peta\byte}), so a literal implementation
would not exist for the environment where the $\Oh(d^2)$ cost matters most.
The claim under test concerns \emph{uplink bytes}, not server FLOPs; the paper
measures ``the number of transmitted parameters with double precision in each
agent''. We therefore compute the identical direction cheaply and report the
true transmission cost. Dropping that metering would make FedNPG-ADMM appear
free by comparison, and the experiment would measure nothing.

\subsection{FedPPO}

The paper names FedPPO and gives its clipping parameter, and specifies nothing further, neither the aggregation scheme nor the local
epoch count. We implement the
FedAvg reading: each agent runs local PPO epochs from the broadcast $\bthe$
and uplinks its parameter delta, and the server averages. The alternative, a
single clipped-surrogate gradient step per round, would parallel
Algorithm~\ref{alg:fednpgadmm}'s structure more neatly, but it strips PPO of
the multi-epoch minibatch reuse that its clipping exists to make safe. Beating
that version would establish nothing. The choice is recorded in every run's
metadata.

\subsection{Communication cost, measured}

\begin{table}[t]
  \centering
  \caption{Measured per-iteration uplink, Swimmer-v4 with four agents
  ($d = \num{4868}$), at 64 bits per parameter. Measured values match the closed
  form exactly.}
  \label{tab:measuredcomm}
  \begin{tabular}{llrr}
    \toprule
    Method & Transmitted & Bytes/iteration & Relative \\
    \midrule
    Standard-average FedNPG & $\bH_i$ ($d^2$) $+\ \bg_i$ & \num{758473344} & $2434\times$ \\
    FedNPG-ADMM             & $\by_i,\ \bg_i$ (2 vectors) & \num{311552}   & $2\times$ \\
    FedPPO                  & $\Delta\bthe_i$ (1 vector)  & \num{155776}   & $1\times$ \\
    \bottomrule
  \end{tabular}
\end{table}

One consequence of Table~\ref{tab:measuredcomm} deserves emphasis, because the
paper's framing invites a reader to assume otherwise: FedNPG-ADMM is not the
cheapest of the three methods. It is $\Oh(d)$, but with twice the constant of
FedPPO, since it transmits two vectors per round where FedPPO transmits one.
The argument for the method therefore rests on what it achieves per round
rather than on what it costs per round, and a figure plotted against iterations
cannot show that. We consequently plot reward against \emph{cumulative bytes},
which is the axis on which the two effects can be read together.

\TODO{insert Figure 3 analogue -- reward vs iterations and reward vs cumulative
bytes, three methods, N=8, three environments, five seeds -- once the 30
baseline runs complete}

\TODO{table of final rewards for the three methods at N=8}

% !TEX root = ../main.tex
\section{Quantizing the agent-to-server transmissions}
\label{sec:quantization}

FedNPG-ADMM reduces communication by shrinking \emph{what} is transmitted from
a matrix to two vectors. A second axis, which the original work does not
explore, is to shrink \emph{how precisely} those vectors are transmitted. The
paper reports overhead at 64 bits per parameter. If performance survives at a
handful of bits per parameter, the saving multiplies the $\Oh(d^2) \to \Oh(d)$
one.

This is the novel contribution of the project. The precedent is Q-GADMM
\citep{elgabli2021qgadmm}, which \citet{lan2023fednpgadmm} cite themselves, and
which quantizes the vectors an ADMM scheme transmits in the \emph{supervised}
setting. We carry that idea into the reinforcement learning setting, where the
transmitted vectors are gradient and consensus directions computed from a
non-stationary sampling distribution. We found no prior work quantizing
agent-to-server transmissions for federated policy gradient with
continuous-control experiments. The nearest results are for compressed temporal
difference learning with linear function approximation
\citep{mitra2024compressedtd}, which is theory-only and not policy gradient.

\subsection{Design}

A quantizer is a \texttt{Transform}: \texttt{tensor} $\rightarrow$
(\texttt{tensor}, \texttt{bits\_per\_element}), attached to the
\texttt{Channel} through a factory that \texttt{train()} accepts. It quantizes
and dequantizes in one call and declares the bit width the wire would have
needed. No file in the algorithm proper knows quantization exists.

Two properties make the experiment measure something. The degraded tensor is
what the server aggregates, because downstream code uses the return value of
\texttt{send} and never the original, so the quantization error enters the
update to $\bthe$. And the declared \texttt{bits\_per\_element} is what the
byte meter uses, so the cost axis moves with the fidelity. Per iteration, with
$R$ ADMM rounds, $N$ agents and $b$ bits, the uplink is
\begin{equation}
  \text{bytes} \;=\; N \,(1 + R)\, d \, b / 8,
  \label{eq:quantbytes}
\end{equation}
since $\bg_i$ is sent once per iteration and $\by_i$ once per round. We use
$R = 1$ throughout, so each agent uplinks $2d$ values per iteration and the
saving relative to the paper's 64-bit convention is exactly $64/b$: $8\times$
at 8 bits and $10.7\times$ at 6. Measured byte counts match
\eqref{eq:quantbytes} exactly.

\paragraph{Stochastic uniform quantization.} For a tensor $x$ at $b$ bits, with
$L = 2^b - 1$ levels and step $s = (\max x - \min x)/L$, the value is placed at
$\lfloor u \rfloor + \mathrm{Bernoulli}(u - \lfloor u \rfloor)$ where
$u = (x - \min x)/s$. Rounding stochastically instead of to nearest makes the
estimator unbiased, $\Exp[\hat{x}] = x$, with per-element variance at most
$s^2/4$. That property is the reason to prefer it here: a deterministic
rounding error is systematic and would accumulate across iterations, whereas an
unbiased one leaves the server receiving the correct gradient and consensus
direction in expectation. The quantizer draws from a dedicated
\texttt{torch.Generator} rather than the global stream, so runs at different
bit widths see identical action-sampling sequences and differ only in the
quantization.

\paragraph{Protocol.} We run the full grid $N \in \{2, 4, 8\}$ $\times$
$b \in \{6, 8\}$ on all three environments, with the same five seeds and the
same hyperparameters as Section~\ref{sec:reproduction} (90 runs). The
unquantized reference is therefore the matching column of
Table~\ref{tab:reproduction}, and every quantized cell differs from it only in
the uplink channel.

\subsection{Results}

\begin{table}[t]
  \centering
  \caption{Final reward under stochastic quantization, mean $\pm$ standard
  deviation across five seeds, computed as in Table~\ref{tab:reproduction}. The
  64-bit rows are the unquantized $N = 2, 4, 8$ columns of
  Table~\ref{tab:reproduction}. Uplink is per agent per iteration, $2db/8$
  bytes, and does not depend on $N$; the federation total is $N$ times it,
  \eqref{eq:quantbytes}.}
  \label{tab:quant}
  \begin{tabular}{llrrrrr}
    \toprule
    Environment & Bits & $N=2$ & $N=4$ & $N=8$ & Uplink (bytes) & Saving \\
    \midrule
    \multirow{3}{*}{Swimmer-v4}
      & 64 & $126.0 \pm 1.3$ & $127.0 \pm 0.8$ & $128.4 \pm 0.8$ & \num{77888} & $1\times$ \\
      & 8  & $126.1 \pm 0.8$ & $127.1 \pm 0.5$ & $128.6 \pm 0.4$ & \num{9736}  & $\mathbf{8.0\times}$ \\
      & 6  & $128.9 \pm 7.4$ & $125.8 \pm 3.4$ & $127.8 \pm 0.7$ & \num{7302}  & $\mathbf{10.7\times}$ \\
    \midrule
    \multirow{3}{*}{Hopper-v4}
      & 64 & $2140 \pm 386$ & $1658 \pm 273$ & $1605 \pm 159$ & \num{295008} & $1\times$ \\
      & 8  & $1901 \pm 593$ & $1755 \pm 302$ & $1713 \pm 150$ & \num{36876}  & $\mathbf{8.0\times}$ \\
      & 6  & $1573 \pm 817$ & $1853 \pm 551$ & $1837 \pm 248$ & \num{27657}  & $\mathbf{10.7\times}$ \\
    \midrule
    \multirow{3}{*}{Humanoid-v4}
      & 64 & $3391 \pm 796$ & $5343 \pm 121$ & $6329 \pm 334$ & \num{11633184} & $1\times$ \\
      & 8  & $2956 \pm 593$ & $5244 \pm 237$ & $5997 \pm 259$ & \num{1454148}  & $\mathbf{8.0\times}$ \\
      & 6  & $497 \pm 34$   & $676 \pm 73$   & $1514 \pm 526$ & \num{1090611}  & $10.7\times$ \\
    \bottomrule
  \end{tabular}
\end{table}

\begin{figure}[t]
  \centering
  \begin{subfigure}[b]{0.49\textwidth}
    %% Return against cumulative uplink traffic at N=8, Swimmer-v4, per bit width.
%% Mean over five seeds; band spans min to max. The 64-bit curve is the N=8
%% series of Figure 2, so the horizontal gap between curves is the saving.
%% Data: figures/quant_comm_swimmer.dat and figures/communication_swimmer.dat
%% (scripts/export_figures.py).
\begin{tikzpicture}
\begin{axis}[
    width=\linewidth,
    height=0.78\linewidth,
    xlabel={Cumulative uplink (bytes)},
    ylabel={Average batch return},
    xmode=log,
    ymin=0, ymax=140,
    ytick distance=35,
    ymajorgrids=true, xmajorgrids=true, grid style=dashed,
    tick label style={font=\footnotesize},
    label style={font=\footnotesize},
    legend style={legend columns=1, font=\footnotesize, cells={anchor=west},
                  anchor=north west, at={(0.02,0.98)}},
]

\addplot[name path=N8qlo, draw=none, forget plot]
    table[x=N8_bytes, y=N8_lo] {figures/communication_swimmer.dat};
\addplot[name path=N8qhi, draw=none, forget plot]
    table[x=N8_bytes, y=N8_hi] {figures/communication_swimmer.dat};
\addplot[colorD, opacity=0.15, forget plot] fill between[of=N8qlo and N8qhi];
\addplot[colorD, line width=0.9pt, solid]
    table[x=N8_bytes, y=N8_mean] {figures/communication_swimmer.dat};
\addlegendentry{64 bit};

\addplot[name path=q8qlo, draw=none, forget plot]
    table[x=q8_bytes, y=q8_lo] {figures/quant_comm_swimmer.dat};
\addplot[name path=q8qhi, draw=none, forget plot]
    table[x=q8_bytes, y=q8_hi] {figures/quant_comm_swimmer.dat};
\addplot[colorA, opacity=0.15, forget plot] fill between[of=q8qlo and q8qhi];
\addplot[colorA, line width=0.9pt, dashed]
    table[x=q8_bytes, y=q8_mean] {figures/quant_comm_swimmer.dat};
\addlegendentry{8 bit};

\addplot[name path=q6qlo, draw=none, forget plot]
    table[x=q6_bytes, y=q6_lo] {figures/quant_comm_swimmer.dat};
\addplot[name path=q6qhi, draw=none, forget plot]
    table[x=q6_bytes, y=q6_hi] {figures/quant_comm_swimmer.dat};
\addplot[colorB, opacity=0.15, forget plot] fill between[of=q6qlo and q6qhi];
\addplot[colorB, line width=0.9pt, dash dot]
    table[x=q6_bytes, y=q6_mean] {figures/quant_comm_swimmer.dat};
\addlegendentry{6 bit};
\end{axis}
\end{tikzpicture}

    \caption{Swimmer-v4}
  \end{subfigure}
  \hfill
  \begin{subfigure}[b]{0.49\textwidth}
    %% Return against cumulative uplink traffic at N=8, Hopper-v4, per bit width.
%% Mean over five seeds; band spans min to max. The 64-bit curve is the N=8
%% series of Figure 2, so the horizontal gap between curves is the saving.
%% Data: figures/quant_comm_hopper.dat and figures/communication_hopper.dat
%% (scripts/export_figures.py).
\begin{tikzpicture}
\begin{axis}[
    width=\linewidth,
    height=0.78\linewidth,
    xlabel={Cumulative uplink (bytes)},
    ylabel={Average batch return},
    xmode=log,
    ymin=0, ymax=3600,
    ytick distance=900,
    scaled y ticks=base 10:-3,
    ymajorgrids=true, xmajorgrids=true, grid style=dashed,
    tick label style={font=\footnotesize},
    label style={font=\footnotesize},
    legend style={legend columns=1, font=\footnotesize, cells={anchor=west},
                  anchor=north west, at={(0.02,0.98)}},
]

\addplot[name path=N8qlo, draw=none, forget plot]
    table[x=N8_bytes, y=N8_lo] {figures/communication_hopper.dat};
\addplot[name path=N8qhi, draw=none, forget plot]
    table[x=N8_bytes, y=N8_hi] {figures/communication_hopper.dat};
\addplot[colorD, opacity=0.15, forget plot] fill between[of=N8qlo and N8qhi];
\addplot[colorD, line width=0.9pt, solid]
    table[x=N8_bytes, y=N8_mean] {figures/communication_hopper.dat};
\addlegendentry{64 bit};

\addplot[name path=q8qlo, draw=none, forget plot]
    table[x=q8_bytes, y=q8_lo] {figures/quant_comm_hopper.dat};
\addplot[name path=q8qhi, draw=none, forget plot]
    table[x=q8_bytes, y=q8_hi] {figures/quant_comm_hopper.dat};
\addplot[colorA, opacity=0.15, forget plot] fill between[of=q8qlo and q8qhi];
\addplot[colorA, line width=0.9pt, dashed]
    table[x=q8_bytes, y=q8_mean] {figures/quant_comm_hopper.dat};
\addlegendentry{8 bit};

\addplot[name path=q6qlo, draw=none, forget plot]
    table[x=q6_bytes, y=q6_lo] {figures/quant_comm_hopper.dat};
\addplot[name path=q6qhi, draw=none, forget plot]
    table[x=q6_bytes, y=q6_hi] {figures/quant_comm_hopper.dat};
\addplot[colorB, opacity=0.15, forget plot] fill between[of=q6qlo and q6qhi];
\addplot[colorB, line width=0.9pt, dash dot]
    table[x=q6_bytes, y=q6_mean] {figures/quant_comm_hopper.dat};
\addlegendentry{6 bit};
\end{axis}
\end{tikzpicture}

    \caption{Hopper-v4}
  \end{subfigure}

  \vspace{0.6em}
  \begin{subfigure}[b]{0.49\textwidth}
    %% Return against cumulative uplink traffic at N=8, Humanoid-v4, per bit width.
%% Mean over five seeds; band spans min to max. The 64-bit curve is the N=8
%% series of Figure 2, so the horizontal gap between curves is the saving.
%% Data: figures/quant_comm_humanoid.dat and figures/communication_humanoid.dat
%% (scripts/export_figures.py).
\begin{tikzpicture}
\begin{axis}[
    width=\linewidth,
    height=0.78\linewidth,
    xlabel={Cumulative uplink (bytes)},
    ylabel={Average batch return},
    xmode=log,
    ymin=0, ymax=7000,
    ytick distance=1750,
    scaled y ticks=base 10:-3,
    ymajorgrids=true, xmajorgrids=true, grid style=dashed,
    tick label style={font=\footnotesize},
    label style={font=\footnotesize},
    legend style={legend columns=1, font=\footnotesize, cells={anchor=west},
                  anchor=north west, at={(0.02,0.98)}},
]

\addplot[name path=N8qlo, draw=none, forget plot]
    table[x=N8_bytes, y=N8_lo] {figures/communication_humanoid.dat};
\addplot[name path=N8qhi, draw=none, forget plot]
    table[x=N8_bytes, y=N8_hi] {figures/communication_humanoid.dat};
\addplot[colorD, opacity=0.15, forget plot] fill between[of=N8qlo and N8qhi];
\addplot[colorD, line width=0.9pt, solid]
    table[x=N8_bytes, y=N8_mean] {figures/communication_humanoid.dat};
\addlegendentry{64 bit};

\addplot[name path=q8qlo, draw=none, forget plot]
    table[x=q8_bytes, y=q8_lo] {figures/quant_comm_humanoid.dat};
\addplot[name path=q8qhi, draw=none, forget plot]
    table[x=q8_bytes, y=q8_hi] {figures/quant_comm_humanoid.dat};
\addplot[colorA, opacity=0.15, forget plot] fill between[of=q8qlo and q8qhi];
\addplot[colorA, line width=0.9pt, dashed]
    table[x=q8_bytes, y=q8_mean] {figures/quant_comm_humanoid.dat};
\addlegendentry{8 bit};

\addplot[name path=q6qlo, draw=none, forget plot]
    table[x=q6_bytes, y=q6_lo] {figures/quant_comm_humanoid.dat};
\addplot[name path=q6qhi, draw=none, forget plot]
    table[x=q6_bytes, y=q6_hi] {figures/quant_comm_humanoid.dat};
\addplot[colorB, opacity=0.15, forget plot] fill between[of=q6qlo and q6qhi];
\addplot[colorB, line width=0.9pt, dash dot]
    table[x=q6_bytes, y=q6_mean] {figures/quant_comm_humanoid.dat};
\addlegendentry{6 bit};
\end{axis}
\end{tikzpicture}

    \caption{Humanoid-v4}
  \end{subfigure}
  \caption{Return against cumulative uplink traffic at $N = 8$, per bit width,
  on the same logarithmic axis as Figure~\ref{fig:communication}. The 64-bit
  curve is the $N = 8$ series of that figure. The solid line is the mean over
  five seeds and the band spans the minimum and maximum seed. Quantization
  translates a curve to the left by $\log_{10}(64/b)$ decades, $0.90$ at 8 bits
  and $1.03$ at 6, without changing its shape whenever the bit width is above
  the threshold: the same return is reached with an order of magnitude less
  traffic. Humanoid at 6 bits is the one curve that drops as well as shifts.}
  \label{fig:quantcomm}
\end{figure}

\begin{figure}[t]
  \centering
  \begin{subfigure}[b]{0.49\textwidth}
    %% Final return against bit width, Swimmer-v4, one series per agent count.
%% Mean +- sd over five seeds (same statistic as Table 1); the 64-bit points are
%% the unquantized N=2,4,8 columns of Table 1. Per-agent uplink is 2d*b/8 bytes,
%% so the x axis is proportional to it. Series are offset by a few percent in x
%% so that the error bars do not overlap.
%% Data: figures/quant_tradeoff_swimmer.dat (scripts/export_figures.py).
\begin{tikzpicture}
\begin{axis}[
    width=\linewidth,
    height=0.78\linewidth,
    xlabel={Bits per transmitted value},
    ylabel={Final return},
    xmode=log, log basis x=2,
    xmin=5, xmax=80,
    xtick={6,8,16,32,64},
    xticklabels={6,8,16,32,64},
    ymin=115, ymax=140,
    ytick distance=5,
    ymajorgrids=true, xmajorgrids=true, grid style=dashed,
    tick label style={font=\footnotesize},
    label style={font=\footnotesize},
    legend style={font=\footnotesize, cells={anchor=west},
                  anchor=south, at={(0.5,0.03)}, legend columns=3},
]

\addplot[colorB, line width=0.9pt, mark=*, mark size=2pt,
        error bars/.cd, y dir=both, y explicit]
    table[x expr=\thisrow{bits}*0.93, y=N2_mean, y error=N2_sd]
    {figures/quant_tradeoff_swimmer.dat};
\addlegendentry{$N=2$};

\addplot[colorC, line width=0.9pt, mark=square*, mark size=2pt,
        error bars/.cd, y dir=both, y explicit]
    table[x expr=\thisrow{bits}*1.0, y=N4_mean, y error=N4_sd]
    {figures/quant_tradeoff_swimmer.dat};
\addlegendentry{$N=4$};

\addplot[colorD, line width=0.9pt, mark=triangle*, mark size=2pt,
        error bars/.cd, y dir=both, y explicit]
    table[x expr=\thisrow{bits}*1.075, y=N8_mean, y error=N8_sd]
    {figures/quant_tradeoff_swimmer.dat};
\addlegendentry{$N=8$};
\end{axis}
\end{tikzpicture}

    \caption{Swimmer-v4}
  \end{subfigure}
  \hfill
  \begin{subfigure}[b]{0.49\textwidth}
    %% Final return against bit width, Hopper-v4, one series per agent count.
%% Mean +- sd over five seeds (same statistic as Table 1); the 64-bit points are
%% the unquantized N=2,4,8 columns of Table 1. Per-agent uplink is 2d*b/8 bytes,
%% so the x axis is proportional to it. Series are offset by a few percent in x
%% so that the error bars do not overlap.
%% Data: figures/quant_tradeoff_hopper.dat (scripts/export_figures.py).
\begin{tikzpicture}
\begin{axis}[
    width=\linewidth,
    height=0.78\linewidth,
    xlabel={Bits per transmitted value},
    ylabel={Final return},
    xmode=log, log basis x=2,
    xmin=5, xmax=80,
    xtick={6,8,16,32,64},
    xticklabels={6,8,16,32,64},
    ymin=0, ymax=3600,
    ytick distance=900,
    scaled y ticks=base 10:-3,
    ymajorgrids=true, xmajorgrids=true, grid style=dashed,
    tick label style={font=\footnotesize},
    label style={font=\footnotesize},
    legend style={font=\footnotesize, cells={anchor=west},
                  anchor=north, at={(0.5,0.97)}, legend columns=3},
]

\addplot[colorB, line width=0.9pt, mark=*, mark size=2pt,
        error bars/.cd, y dir=both, y explicit]
    table[x expr=\thisrow{bits}*0.93, y=N2_mean, y error=N2_sd]
    {figures/quant_tradeoff_hopper.dat};
\addlegendentry{$N=2$};

\addplot[colorC, line width=0.9pt, mark=square*, mark size=2pt,
        error bars/.cd, y dir=both, y explicit]
    table[x expr=\thisrow{bits}*1.0, y=N4_mean, y error=N4_sd]
    {figures/quant_tradeoff_hopper.dat};
\addlegendentry{$N=4$};

\addplot[colorD, line width=0.9pt, mark=triangle*, mark size=2pt,
        error bars/.cd, y dir=both, y explicit]
    table[x expr=\thisrow{bits}*1.075, y=N8_mean, y error=N8_sd]
    {figures/quant_tradeoff_hopper.dat};
\addlegendentry{$N=8$};
\end{axis}
\end{tikzpicture}

    \caption{Hopper-v4}
  \end{subfigure}

  \vspace{0.6em}
  \begin{subfigure}[b]{0.49\textwidth}
    %% Final return against bit width, Humanoid-v4, one series per agent count.
%% Mean +- sd over five seeds (same statistic as Table 1); the 64-bit points are
%% the unquantized N=2,4,8 columns of Table 1. Per-agent uplink is 2d*b/8 bytes,
%% so the x axis is proportional to it. Series are offset by a few percent in x
%% so that the error bars do not overlap.
%% Data: figures/quant_tradeoff_humanoid.dat (scripts/export_figures.py).
\begin{tikzpicture}
\begin{axis}[
    width=\linewidth,
    height=0.78\linewidth,
    xlabel={Bits per transmitted value},
    ylabel={Final return},
    xmode=log, log basis x=2,
    xmin=5, xmax=80,
    xtick={6,8,16,32,64},
    xticklabels={6,8,16,32,64},
    ymin=0, ymax=7000,
    ytick distance=1750,
    scaled y ticks=base 10:-3,
    ymajorgrids=true, xmajorgrids=true, grid style=dashed,
    tick label style={font=\footnotesize},
    label style={font=\footnotesize},
    legend style={font=\footnotesize, cells={anchor=west},
                  anchor=south east, at={(0.97,0.03)}, legend columns=1},
]

\addplot[colorB, line width=0.9pt, mark=*, mark size=2pt,
        error bars/.cd, y dir=both, y explicit]
    table[x expr=\thisrow{bits}*0.93, y=N2_mean, y error=N2_sd]
    {figures/quant_tradeoff_humanoid.dat};
\addlegendentry{$N=2$};

\addplot[colorC, line width=0.9pt, mark=square*, mark size=2pt,
        error bars/.cd, y dir=both, y explicit]
    table[x expr=\thisrow{bits}*1.0, y=N4_mean, y error=N4_sd]
    {figures/quant_tradeoff_humanoid.dat};
\addlegendentry{$N=4$};

\addplot[colorD, line width=0.9pt, mark=triangle*, mark size=2pt,
        error bars/.cd, y dir=both, y explicit]
    table[x expr=\thisrow{bits}*1.075, y=N8_mean, y error=N8_sd]
    {figures/quant_tradeoff_humanoid.dat};
\addlegendentry{$N=8$};
\end{axis}
\end{tikzpicture}

    \caption{Humanoid-v4}
  \end{subfigure}
  \caption{Final return against bit width for each agent count, mean $\pm$
  standard deviation across five seeds (Table~\ref{tab:quant}). Per-agent uplink
  is $2db/8$ bytes, so the horizontal axis is proportional to it, and moving
  from the right-hand points to the left-hand ones divides the traffic by $8$
  and $10.7$. Series are offset slightly in $x$ so the error bars stay
  readable. The Swimmer axis is zoomed to the range of the data. A flat line
  is the desired outcome; only Humanoid at 6 bits breaks it.}
  \label{fig:quanttradeoff}
\end{figure}

\paragraph{8 bits cost nothing measurable.} At 8 bits the uplink falls by
$8\times$ on every environment, and Table~\ref{tab:quant} shows no cell in
which the reward pays for it. On Swimmer the quantized means sit within $0.2$
of the unquantized ones at all three agent counts, with a seed spread that is,
if anything, tighter ($\pm 0.4$ against $\pm 0.8$ at $N=8$). On Humanoid, the
environment with the most at stake, the gap is $99$ at $N=4$ and $332$ at
$N=8$, each inside one standard deviation of the reference, and the $N=2$ gap
of $435$ is inside the $\pm 796$ spread of the unquantized runs. On Hopper the
8-bit means are higher than the reference at $N = 4$ and $8$ and lower at
$N = 2$, all within the seed spread. Figure~\ref{fig:quanttradeoff} draws the
same result as flat lines, and Figure~\ref{fig:quantcomm} shows it along the
whole run rather than only at the end: the 8-bit curve has the shape of the
64-bit one, displaced by almost a decade to the left.

In absolute terms the saving is large. A full $N=8$ Humanoid run uplinks
\SI{279}{\giga\byte} at 64 bits and \SI{34.9}{\giga\byte} at 8, so 8-bit
quantization removes \SI{244}{\giga\byte} of agent-to-server traffic from a
single training run at no measurable cost in return. The corresponding figures
are \SI{4.72}{\giga\byte} to \SI{0.59}{\giga\byte} on Hopper and
\SI{623}{\mega\byte} to \SI{78}{\mega\byte} on Swimmer.

\paragraph{The saving compounds with the $\Oh(d)$ one.} Standard-average FedNPG
uplinks $d^2 + d$ values per agent per round at 64 bits; 8-bit FedNPG-ADMM
uplinks $2d$ values at 8 bits. The ratio is $4(d+1)$, which turns the
$2434\times$, $9219\times$ and $363537\times$ reductions of
Section~\ref{sec:baselines} into $1.9 \times 10^4$, $7.4 \times 10^4$ and
$2.9 \times 10^6$ on Swimmer, Hopper and Humanoid. Quantization also changes
the ordering in Table~\ref{tab:measuredcomm}. At 64 bits FedNPG-ADMM transmits
twice what FedPPO does; at 8 bits it transmits $16d$ bits per agent per
iteration against FedPPO's $64d$, a quarter of the unquantized first-order
baseline. FedPPO's parameter delta could in principle be quantized as well,
which we did not test, so this is a statement about the configurations we ran
and not a ranking of the two methods under equal compression.

The converse does not hold: quantization cannot rescue standard-average
FedNPG. To match 8-bit FedNPG-ADMM, its $d^2 + d$ values would have to travel
at $16/(d+1)$ bits each, $3.3 \times 10^{-3}$ bits on Swimmer and less on the
larger policies. No quantizer reaches below one bit per value, so the
$\Oh(d^2) \to \Oh(d)$ reformulation is a precondition for the precision axis
to matter, not an alternative to it.

\paragraph{6 bits: a threshold that depends on $d$.} At 6 bits the saving
rises to $10.7\times$. Swimmer ($d = \num{4868}$) and Hopper
($d = \num{18438}$) keep their returns, while Humanoid ($d = \num{727074}$)
collapses at every agent count, to between $13\%$ and $24\%$ of the
unquantized return. Two qualifications apply to the surviving environments. On
Swimmer the 6-bit spread widens at $N = 2$ and $N = 4$ ($\pm 7.4$ and $\pm 3.4$
against $\pm 1.3$ and $\pm 0.8$). The $N=2$ mean of $128.9$ is pulled up by a
single seed at $142.0$, with the other four between $124.9$ and $126.8$, and at
$N=4$ two seeds end near $123$; this is the first visible effect of the
quantization noise, and it is gone at $N = 8$. On Hopper the 6-bit means show no
ordering in bit width ($1573$, $1853$, $1837$ against $2140$, $1658$, $1605$),
and the differences are inside the seed spread that Section~\ref{sec:hopper}
documents for this environment. The Hopper column measures seed variance, not
quantization.

That the threshold moves with $d$ is consistent with how the quantizer sets
its step. The step $s$ is fixed by the range of the whole tensor, and in a
vector of $\num{727074}$ entries the extremes sit much further out in the
tails than in one of $\num{4868}$. Most entries are then small compared with
$s$, and what reaches the server for them is dominated by rounding noise,
correct only in expectation. We have not measured the per-tensor ranges, so this reading is a
hypothesis; a per-bit sweep on a fourth environment of intermediate $d$ would
test it.

The collapsed Humanoid runs are at least consistent with the reason we
chose an unbiased quantizer. At 6 bits, return grows $3.0\times$ from $N=2$ to
$N=8$ ($497 \to 1514$), against $1.9\times$ for the unquantized runs
($3391 \to 6329$), so the fraction of the reference retained rises from
$15\%$ and $13\%$ at $N = 2$ and $4$ to $24\%$ at $N=8$. Each agent's
quantization error is independent and zero-mean, so the server's average over
$N$ agents divides its variance by $N$, and more agents buy back part of the
fidelity the coarse quantizer took away. A biased quantizer would leave the
same systematic error after averaging. With three points and a non-monotone
middle one, this is suggestive rather than conclusive.

\subsection{What the byte counts omit}

The reported bytes are $d \cdot b / 8$ per transmission and exclude the
quantization metadata a real deployment would also send: $\min x$ and the
scale, two \texttt{float32} values or 64 bits per transmission. Relative to a
$d \cdot b$-bit payload the overhead is $64/(d\,b)$. At the bit widths we run
this is at most $0.22\%$ (Swimmer, 6 bits), $0.06\%$ on Hopper and below
$0.002\%$ on Humanoid, so the $8\times$ saving is $7.99\times$ once the
metadata is paid for. The downlink broadcast of $\bthe$ and $\by$ is not
counted either, for the same reason it is not counted in the paper. Neither
omission changes the ordering in Table~\ref{tab:quant}, but both make the
reported compression slightly optimistic.

\paragraph{Skipped steps.} Quantization noise can in principle drive
$(\sum_i \bg_i)^{\top}\by$ negative, leaving line~12's step size undefined
(Section~\ref{sec:implementation}). On Hopper and Humanoid this never occurs
in any of the 60 quantized runs. On Swimmer it does, but it does not track the
bit width: the fifteen 8-bit runs skip 51 of their \num{15000} iterations in
total (at most 12 in one run), while the fifteen 6-bit runs skip 3. Coarser
quantization does not make the step more often undefined.

% !TEX root = ../main.tex
\section{Discussion}
\label{sec:discussion}

\subsection{Hypotheses tested and refuted}

Hopper's inverted scaling (Section~\ref{sec:hopper}) absorbed the larger part
of our experimental effort. We record below the hypotheses that measurement
subsequently refuted. They are worth setting out rather than omitting, both
because each one narrows the space in which the remaining explanation must sit,
and because the order in which they failed is informative about where the
intuitions behind this method mislead.

\begin{description}
  \item[ADMM consensus quality.] At $N=8$ the agents' $\by_i$ are mutually
    near-orthogonal (mean pairwise cosine $\approx 0.05$), and their average
    retains only $\approx 0.57$ alignment with the true pooled natural-gradient
    direction against $0.75$--$0.83$ for a single agent. Forcing consensus with five ADMM rounds per policy iteration raised the
    cosine to $\approx 0.76$ and made performance \emph{worse}, and an exact
    pooled solve, which has zero consensus error by construction, landed within
    noise of one round. Consensus quality is not the bottleneck.
  \item[Trust-region overshoot.] Line~12's closed-form step is exact only at
    ADMM convergence, so with one round it could in principle overshoot. The
    measured overshoot factor was $0.68$--$1.04$ and the achieved KL
    $0.0054$--$0.0107$ against $\delta = 0.01$, so the step is mildly
    \emph{conservative} at both $N=1$ and $N=8$.
  \item[Architecture.] Table~4 mandates ReLU hidden layers and a Tanh-squashed
    policy mean, whereas no reference implementation squashes the mean. Testing
    all four combinations, the Table~4 configuration won, and removing the Tanh
    alone was the worst. The natural gradient largely absorbs the
    reparameterization, so the apparent gradient attenuation does not bite.
  \item[Exploration noise.] Three interventions left the fall rate at
    $9$--$12\%$ (Section~\ref{sec:variance}): releasing the log-std floor,
    releasing the ceiling, and annealing the noise to a level at which a fixed
    policy never falls.
\end{description}

The explanation that survives is the reward-structure argument of
Section~\ref{sec:hopper}, which is consistent with the known tendency of
MuJoCo survival bonuses to create local optima \citep{mania2018ars,
reda2020locomote}.

\subsection{Two observations we could not find in the literature}

We located no published source for either of the following, and present them
as our own measurements. First, a quantified failure rate and bimodal return
distribution for a trained stochastic MuJoCo policy. The literature states
that the episode ends when the robot falls, but not how often that happens to
a converged policy. Second, per-dimension divergence of a Gaussian policy's
$\log\sigma$: with the floor released, one action dimension collapsed to
$0.056$ while another drifted up to $0.311$, leaving the mean unchanged. TRPO
has no entropy term and the trust region bounds the KL of the joint
distribution, so a dimension whose actions barely affect the advantage
receives almost no gradient on its $\log\sigma$.

\subsection{Threats to validity}

Our study uses five seeds against the original's ten, so our error bars are
wider than they appear relative to a ten-seed estimate. All runs use a single
machine and a single software stack, so implementation-level differences from
the authors' setup cannot be separated from algorithmic ones. That concern has
precedent here: the same algorithm under identical hyperparameters has been
shown to differ by more than a factor of two on Hopper across implementations
\citep{henderson2018matters}. Finally, two hyperparameters ($\eta$ and the
interpretation of Table~4's learning-rate column,
Section~\ref{sec:implementation}) are our reading of an ambiguous source, and
a different reading could change the Hopper conclusion in particular.

% !TEX root = ../main.tex
\section{Conclusion}
\label{sec:conclusion}

We reproduced FedNPG-ADMM from scratch and evaluated it across three MuJoCo
environments, four agent counts and five seeds. Swimmer reproduces the
reported final rewards at every agent count, monotone in $N$ and with tighter
seed variance than the original throughout, and Humanoid scales
$6.4\times$ from one agent to eight. Hopper does not reproduce. Our
single-agent result exceeds the published value, performance then decreases
with more agents, and we attribute this to a speed-versus-survival local
optimum that Hopper's small per-step survival bonus makes attractive and that
reduced gradient noise makes easier to fall into.

We measured the communication argument directly instead of assuming it,
obtaining per-agent per-round reductions of $2434\times$, $9219\times$ and
$363537\times$ on the three environments. FedNPG-ADMM is not the cheapest of
the compared methods in absolute terms, since it transmits twice what FedPPO
does, so its case rests on convergence per round.

Quantizing the agent-to-server transmissions, which the original work does not
attempt, removes most of what remains of that cost. At 8 bits the uplink falls
by $8\times$ on every environment and agent count, with final rewards inside
one standard deviation of the unquantized runs over five seeds; on Humanoid
that is \SI{244}{\giga\byte} less traffic per $N=8$ run. Compounded with the
$\Oh(d^2) \to \Oh(d)$ reduction it reaches $1.9\times 10^{4}$,
$7.4\times 10^{4}$ and $2.9\times 10^{6}$ against standard FedNPG, and it
brings FedNPG-ADMM's uplink to a quarter of unquantized FedPPO's. At 6 bits
($10.7\times$) Swimmer and Hopper are still unaffected while Humanoid
collapses at every agent count, so the safe bit width shrinks as the policy
grows.

\paragraph{Limitations and future work.} The missing axis is communication
\emph{frequency}. FedNPG-ADMM optimizes what is transmitted per round, while
FedAvg-style schemes with $E$ local steps optimize how often rounds occur
\citep{jin2022fedrlheterogeneity, khodadadian2022linearspeedup}. The two are
orthogonal and could be combined. Our Hopper diagnosis raises a second
question we did not pursue: whether the exploration lost to variance reduction
can be restored through an entropy term, which Algorithm~\ref{alg:fednpgadmm}
lacks, without giving up the convergence benefit that more agents provide.


\bibliographystyle{plainnat}
\bibliography{refs}

\appendix
% !TEX root = ../main.tex
\section{Hyperparameters}
\label{app:hyperparameters}

\begin{table}[t]
  \centering
  \caption{Values transcribed from Table~4 of \citet{lan2023fednpgadmm}.
  $\lambda_{\text{GAE}} = 0.95$ comes from the paper's setup text rather than
  the table.}
  \label{tab:table4}
  \begin{tabular}{lrrr}
    \toprule
     & Swimmer-v4 & Hopper-v4 & Humanoid-v4 \\
    \midrule
    ADMM penalty $\rho$        & $0.1$   & $0.1$   & $0.01$ \\
    Trust region $\delta$      & $0.01$  & $0.01$  & $0.01$ \\
    Discount $\gamma$          & $0.99$  & $0.99$  & $0.99$ \\
    Timesteps per iteration $T$& $2048$  & $1024$  & $512$ \\
    Iterations $K$             & $1000$  & $2000$  & $3000$ \\
    Learning rate (critic)     & $3\times10^{-4}$ & $3\times10^{-4}$ & $10^{-5}$ \\
    Hidden layers              & $64\times64$ & $128\times128$ & $512^{3}$ \\
    Hidden activation          & ReLU & ReLU & ReLU \\
    Policy mean output         & Tanh & Tanh & Tanh \\
    \bottomrule
  \end{tabular}
\end{table}

\begin{table}[t]
  \centering
  \caption{Settings not specified by the paper. Each follows the practical
  basis the paper's setup section cites: TRPO with line search, and
  stable-baselines3.}
  \label{tab:ours}
  \begin{tabular}{lll}
    \toprule
    Setting & Value & Basis \\
    \midrule
    Observation normalization & running mean/std, clip $10$ & SB3 \texttt{VecNormalize} \\
    Reward normalization      & scale-only by running return std, clip $10$ & SB3 \texttt{VecNormalize} \\
    Advantage normalization   & per agent, $(a-\mu)/(\sigma+10^{-8})$ & SB3 \\
    Orthogonal initialization & $\sqrt{2}$ hidden, $0.01$ policy head, $1.0$ value head & SB3 \\
    $\log\sigma$ initial      & $-0.5$ & tianshou, Spinning Up \\
    $\log\sigma$ clamp        & $[-2.0,\, 0.5]$ & ours \\
    Critic fitting            & 10 epochs, minibatch 64, Adam & SB3 PPO defaults \\
    CG iterations / damping   & $10$ / $\rho$ & classic TRPO; damping $=\rho$ per line~8 \\
    Line search               & 10 backtracks, shrink $0.8$ & TRPO \\
    $\eta$                    & $1.0$ & not specified anywhere \\
    Seeds                     & $0$--$4$ & --- \\
    \bottomrule
  \end{tabular}
\end{table}


\end{document}
